\documentclass[12pt,a4paper]{article}

\usepackage[margin=2.5cm]{geometry}
\usepackage{hyperref}
\usepackage{listings}
\usepackage{xcolor}
\usepackage{parskip}
\usepackage{titlesec}
\usepackage{enumitem}
\usepackage{booktabs}
\usepackage{array}
\usepackage{fancyhdr}
\usepackage{fontenc}
\usepackage{inputenc}

% Colour scheme
\definecolor{codeblue}{RGB}{30,100,180}
\definecolor{codegray}{RGB}{245,245,245}
\definecolor{codecomment}{RGB}{100,100,100}
\definecolor{codestring}{RGB}{180,60,40}
\definecolor{accentblue}{RGB}{0,80,160}

% Code listing style
\lstdefinestyle{codestyle}{
    backgroundcolor=\color{codegray},
    commentstyle=\color{codecomment}\itshape,
    keywordstyle=\color{codeblue}\bfseries,
    stringstyle=\color{codestring},
    basicstyle=\ttfamily\small,
    breakatwhitespace=false,
    breaklines=true,
    captionpos=b,
    keepspaces=true,
    showspaces=false,
    showstringspaces=false,
    showtabs=false,
    tabsize=2,
    frame=single,
    framerule=0.4pt,
    rulecolor=\color{gray!40}
}

\lstset{style=codestyle}

% Section formatting
\titleformat{\section}{\large\bfseries\color{accentblue}}{}{0em}{}[\titlerule]
\titleformat{\subsection}{\normalsize\bfseries}{}{0em}{}

% Header/footer
\pagestyle{fancy}
\fancyhf{}
\rhead{\textcolor{gray}{\small Sheffield Junior Programmer Exercise}}
\lhead{\textcolor{gray}{\small README File}}
\cfoot{\thepage}
\renewcommand{\headrulewidth}{0.4pt}

\hypersetup{
    colorlinks=true,
    linkcolor=accentblue,
    urlcolor=codeblue
}

% ─────────────────────────────────────────────────────────────────────────────
\begin{document}

\begin{titlepage}
    \centering
    \vspace*{3cm}
    {\huge\bfseries\color{accentblue} Sheffield Junior\\[0.4em] Programmer Exercise\\[0.4em]}
    {\large README File\\[2.5cm]}
    \rule{\linewidth}{0.5pt}\\[0.6em]
    {\large\itshape A Python application demonstrating data ingestion,\\
    a REST API, and a CSV export pipeline backed by MySQL}\\[0.6em]
    \rule{\linewidth}{0.5pt}
    \vfill
    {\normalsize March 2026}
\end{titlepage}

\tableofcontents
\newpage

% ─────────────────────────────────────────────────────────────────────────────
\section{Running the Application}

This section provides complete, step-by-step instructions for setting up and
running every component of the application from a clean checkout.

\subsection{Prerequisites}

Before proceeding, ensure the following software is installed and available on
your \texttt{PATH}:

\begin{itemize}[leftmargin=2em]
    \item \textbf{Python 3.11+} --- the application uses f-string features and
          walrus operators that require Python 3.8 at minimum; 3.11 is
          recommended.
    \item \textbf{Docker} --- the database setup script automatically pulls and
          starts a \texttt{mysql:latest} container. Docker must be running
          before any Python script is executed.
    \item \textbf{pip} --- used to install Python dependencies.
\end{itemize}

\subsection{Clone the Repository and Install Dependencies}

\begin{lstlisting}[language=bash, caption={Install Python dependencies}]
# From the root of the repository
pip install -r requirements.txt
\end{lstlisting}

All required packages are pinned in \texttt{requirements.txt}. Key runtime
dependencies include \texttt{mysql-connector-python}, \texttt{Flask},
\texttt{python-dotenv}, and \texttt{rich}.

\subsection{Configure the Environment File}

The application reads all configuration from a \texttt{.env} file in the
project root. An example is provided:

\begin{lstlisting}[language=bash, caption={Create your \texttt{.env} file}]
cp .env.example .env
\end{lstlisting}

The default values work out of the box if you are running Docker locally.
Edit \texttt{.env} only if you need a different port or database password:

\begin{center}
\begin{tabular}{ll}
\toprule
\textbf{Variable} & \textbf{Default} \\
\midrule
\texttt{DB\_HOST}     & \texttt{localhost} \\
\texttt{DB\_NAME}     & \texttt{shopdb}    \\
\texttt{DB\_USER}     & \texttt{root}      \\
\texttt{DB\_PASSWORD} & \texttt{p@ssw0rd}  \\
\texttt{DB\_PORT}     & \texttt{3306}      \\
\texttt{API\_PORT}    & \texttt{8000}      \\
\texttt{OUTPUT\_DIR}  & \texttt{output}    \\
\bottomrule
\end{tabular}
\end{center}

\subsection{Step 1 --- Set Up the Database and Import Data}

Run \texttt{task\_one.py} to seed the database. This script will:

\begin{enumerate}[leftmargin=2em]
    \item Invoke \texttt{utilities/dbsetup.py}, which checks whether a Docker
          container named \texttt{mySQL-shop} is already running. If not, it
          either restarts an existing stopped container or creates a fresh one
          with \texttt{docker run}.
    \item Wait (with a live progress bar) until MySQL is accepting connections.
    \item Create the \texttt{shopdb} database if it does not exist.
    \item Read \texttt{data/customers\_data.csv} and
          \texttt{data/orders\_data.csv}, creating the \texttt{customers} and
          \texttt{orders} tables automatically, then inserting every row that
          does not already exist.
\end{enumerate}

\begin{lstlisting}[language=bash, caption={Populate the database}]
python task_one.py
\end{lstlisting}

Expected output:
\begin{lstlisting}[caption={Successful run output}]
Task 1 has run successfully.
The database has been populated with data from the CSV files.
\end{lstlisting}

\subsection{Step 2 --- Start the REST API}

\begin{lstlisting}[language=bash, caption={Start the Flask API server}]
python task_two.py
\end{lstlisting}

The Flask development server starts on the port defined by \texttt{API\_PORT}
(default \texttt{8000}). Keep this terminal open while using the API.

\textbf{Available endpoint:}

\begin{center}
\begin{tabular}{>{\ttfamily}l l}
\toprule
\textbf{URL} & \textbf{Description} \\
\midrule
GET /customer/\{id\}            & Returns customer profile and orders as JSON \\
GET /customer/\{id\}?pretty=true & Returns the same data formatted as HTML \\
\bottomrule
\end{tabular}
\end{center}

Example requests:

\begin{lstlisting}[language=bash, caption={Example API calls}]
# JSON response
curl http://localhost:8000/customer/CUS2279

# HTML response
curl http://localhost:8000/customer/CUS2279?pretty=true

# Open in a browser for the HTML view
open http://localhost:8000/customer/CUS2279?pretty=true
\end{lstlisting}

\subsection{Step 3 --- Run the Export Script}

\begin{lstlisting}[language=bash, caption={Generate CSV exports}]
python task_three.py
\end{lstlisting}

This produces two files in the directory specified by \texttt{OUTPUT\_DIR}
(default \texttt{output/}):

\begin{itemize}[leftmargin=2em]
    \item \texttt{output/active\_customers.csv} --- all customers whose
          \texttt{Status} is \texttt{Active}, with a computed \texttt{Full Name}
          column appended.
    \item \texttt{output/active\_customer\_orders.csv} --- all orders belonging
          to those active customers, with a computed \texttt{Total Price} column
          appended.
\end{itemize}

Reference outputs are available in \texttt{output.example/} for comparison.

\subsection{Optional --- Regenerate Sample Data}

If you wish to produce a fresh randomised dataset before running the pipeline:

\begin{lstlisting}[language=bash, caption={Regenerate sample CSV data}]
python utilities/data_generator.py
\end{lstlisting}

This overwrites \texttt{data/customers\_data.csv} and
\texttt{data/orders\_data.csv} with 50 new synthetic records each.

% ─────────────────────────────────────────────────────────────────────────────
\section{Design \& Implementation}

\subsection{Database: MySQL via Docker}

MySQL was chosen as the relational database for several practical reasons:

\begin{itemize}[leftmargin=2em]
    \item \textbf{Structured, relational data} --- the dataset consists of two
          clearly related entities (customers and orders) with a foreign-key
          relationship on \texttt{CustomerID}. A relational database is a
          natural and well-supported fit.
    \item \textbf{Ubiquity and tooling} --- MySQL is one of the most widely
          deployed databases; it has mature Python drivers, excellent
          documentation, and broad familiarity.
    \item \textbf{Docker for portability} --- wrapping MySQL in a Docker
          container removes the need for reviewers to install and configure a
          local MySQL server. The setup script in \texttt{utilities/dbsetup.py}
          handles container lifecycle automatically (start, restart, create),
          making the application self-contained with zero manual DB
          configuration.
    \item \textbf{Schema-free table creation} --- because the CSV column
          names are used directly to create the table schema at insertion time,
          the application can adapt to different CSV structures without
          hard-coded \texttt{CREATE TABLE} statements. All columns use
          \texttt{VARCHAR(255)}, which is flexible enough for the data at hand.
\end{itemize}

A lightweight alternative like SQLite would have eliminated the Docker
dependency entirely, but MySQL better reflects a production-like setup and
demonstrates networking, container orchestration, and real TCP connectivity.

\subsection{Web Framework: Flask}

Flask was selected for the API layer because:

\begin{itemize}[leftmargin=2em]
    \item \textbf{Minimal overhead} --- the API surface is a single endpoint.
          Flask imposes no boilerplate beyond what is needed, keeping the code
          concise.
    \item \textbf{Flexibility} --- the custom \texttt{optionalParam} decorator
          in \texttt{utilities/api.py} integrates cleanly with Flask's request
          context, allowing query-string parameters to be declared and type-
          validated inline with the view function.
    \item \textbf{Dual-format responses} --- the same endpoint serves both JSON
          (for machine consumers) and HTML (for browser viewing) via the
          \texttt{?pretty=true} query parameter. Flask's minimal structure made
          it straightforward to branch on this parameter and return either
          \texttt{dict} (auto-serialised to JSON by Flask) or a raw HTML string.
\end{itemize}

\subsection{Key Libraries}

\begin{center}
\begin{tabular}{>{\ttfamily}l p{10cm}}
\toprule
\textbf{Library} & \textbf{Reason for inclusion} \\
\midrule
mysql-connector-python & Official MySQL driver for Python; provides the
                          \texttt{cursor} interface used throughout
                          \texttt{database.py}. \\[0.4em]
python-dotenv / dotenv & Reads the \texttt{.env} file; cleanly separates
                          secrets and config from source code, following
                          twelve-factor app principles. \\[0.4em]
rich                   & Renders the animated spinner and progress bar while
                          waiting for MySQL to become available, giving clear
                          visual feedback during startup. \\[0.4em]
flask-marshmallow /    & Included in \texttt{requirements.txt} for potential
Flask-SQLAlchemy         structured serialisation and ORM use; not yet
                          actively used in the current codebase. \\[0.4em]
names / rstr           & Used by \texttt{data\_generator.py} to produce
                          realistic synthetic customer records with valid-
                          looking names, emails, and UK phone numbers. \\
\bottomrule
\end{tabular}
\end{center}

\subsection{Notable Design Decisions}

\textbf{SQL injection guard.} Rather than switching to parameterised queries
at the driver level (the conventional approach), a decorator \texttt{sql\_guard}
was implemented to check all string arguments for injection tokens
(\texttt{;} and \texttt{--}) before any SQL is executed. This is a defence-in-
depth measure layered on top of the existing query structure.

\textbf{Table and schema auto-creation.} The \texttt{Database.ensureTable}
method creates the table on the first insert, deriving column names from the
CSV headers. This avoids any migration tooling for a project of this scale and
keeps setup to a single script run.

\textbf{Dual-format API.} The decision to return both JSON and HTML from the
same endpoint --- controlled by \texttt{?pretty=true} --- was made to support
both programmatic API consumers and quick human inspection in a browser without
requiring a separate frontend server.

\textbf{Functional column transformation in \texttt{task\_three.py}.} The
\texttt{addNewField} function is written in a generic, functional style: it
accepts a list of field names and an arbitrary \texttt{operation} callable,
then appends the result as a new column. This means the same function handles
both the string concatenation of \texttt{Full Name} and the numeric
multiplication of \texttt{Total Price}, keeping the export logic DRY.

% ─────────────────────────────────────────────────────────────────────────────
\section{Application Flow}

This section traces data through all three tasks and the supporting utilities.

\subsection{Overview}

\begin{center}
\begin{tabular}{ccccc}
\textbf{CSV files} & $\rightarrow$ & \textbf{task\_one.py} & $\rightarrow$ & \textbf{MySQL (Docker)} \\[0.4em]
\textbf{MySQL} & $\rightarrow$ & \textbf{task\_two.py} & $\rightarrow$ & \textbf{HTTP endpoint} \\[0.4em]
\textbf{MySQL} & $\rightarrow$ & \textbf{task\_three.py} & $\rightarrow$ & \textbf{Output CSVs} \\
\end{tabular}
\end{center}

\subsection{Task 1 --- Data Ingestion}

\begin{enumerate}[leftmargin=2em]
    \item \texttt{task\_one.py} imports \texttt{Database} from
          \texttt{utilities/database.py}. This import triggers the module-level
          setup code in \texttt{database.py}: \texttt{dbsetup.run()} is called,
          which inspects running Docker containers and either creates, restarts,
          or does nothing with the \texttt{mySQL-shop} container.
          \texttt{dbsetup.wait()} then polls MySQL on the configured port until
          it accepts connections (up to 30 attempts, 2 seconds apart).

    \item A MySQL connection is established and a buffered cursor created.
          The \texttt{shopdb} database is created if absent, then selected.

    \item For each CSV file (\texttt{customers\_data.csv},
          \texttt{orders\_data.csv}), the script:
          \begin{itemize}
              \item Opens the file and reads the header row, stripping spaces
                    from column names (e.g. \texttt{First Name} becomes
                    \texttt{FirstName}).
              \item For each subsequent row, calls
                    \texttt{Database.getRow(table, primaryKey, id)} to check
                    whether a record with that ID already exists.
              \item If no existing record is found, calls
                    \texttt{Database.insert(table, columns, row)}, which first
                    ensures the table exists (creating it with all-\texttt{VARCHAR(255)}
                    columns if not), then executes the \texttt{INSERT} statement.
          \end{itemize}

    \item Every insert is immediately committed. The idempotency check
          (\texttt{getRow} before \texttt{insert}) means re-running
          \texttt{task\_one.py} is safe and will not duplicate rows.
\end{enumerate}

\subsection{Task 2 --- REST API}

\begin{enumerate}[leftmargin=2em]
    \item \texttt{task\_two.py} imports \texttt{Database} (triggering the same
          Docker/MySQL startup described above) and registers a single Flask
          route: \texttt{GET /customer/<id>}.

    \item The view function \texttt{fetchCustomer} is decorated with
          \texttt{@api.optionalParam("pretty", False, type=bool)}.
          At request time the decorator reads \texttt{?pretty} from the query
          string, coerces it to \texttt{bool}, falls back to \texttt{False} if
          the coercion fails, and sets it as an attribute on the wrapper
          function so the body can access it as \texttt{fetchCustomer.pretty}.

    \item \textbf{JSON path} (\texttt{pretty=False}, the default):
          \begin{itemize}
              \item \texttt{Database.getRow("customers", "UniqueID", id)} is
                    called. If no record is returned, a \texttt{\{"error": "Customer not found"\}}
                    dict is returned, which Flask serialises to JSON with a
                    \texttt{200} status.
              \item Otherwise, all orders for that customer are fetched from
                    \texttt{Database.getRow("orders", "CustomerID", customer[0])}.
              \item A structured response dict is assembled with a
                    \texttt{"Customer Profile"} key and, if orders exist, an
                    \texttt{"Orders"} key whose sub-keys are order IDs. Flask
                    auto-serialises this to JSON.
          \end{itemize}

    \item \textbf{HTML path} (\texttt{pretty=true}):
          \begin{itemize}
              \item The customer tuple is passed through
                    \texttt{utilities/html.py}'s \texttt{htmlClean} function,
                    which applies \texttt{html.escape} to every field, guarding
                    against stored XSS before the values are interpolated into
                    an HTML template string.
              \item Customer details are rendered as an HTML table. Orders, if
                    present, are appended as additional tables below.
              \item The raw HTML string is returned directly; Flask sends it
                    with a \texttt{text/html} content type.
          \end{itemize}

    \item \texttt{api.start()} calls \texttt{app.run(port=env("API\_PORT"))},
          starting the development server.
\end{enumerate}

\subsection{Task 3 --- Export Script}

\begin{enumerate}[leftmargin=2em]
    \item All column names for the \texttt{customers} table are retrieved via
          \texttt{Database.getHeaders("customers")} (which issues
          \texttt{SHOW COLUMNS}).

    \item Active customers are fetched with
          \texttt{Database.getRow("customers", "Status", "Active")}, returning
          a list of tuples.

    \item \texttt{addNewField} is called to append a \texttt{Full Name} column:
          it iterates over every customer tuple, calls \texttt{wordjoin(FirstName, Surname)}
          (a simple \texttt{" ".join(*args)}), and appends the result to the
          tuple. The \texttt{customer\_headers} list is also extended in-place.

    \item For each active customer a dictionary entry is created mapping their
          \texttt{UniqueID} to their orders (fetched by
          \texttt{CustomerID}).

    \item Order numeric fields (\texttt{UnitPrice}, \texttt{Quantity}) are
          converted from strings to Python numbers using \texttt{eval}. A
          \texttt{Total Price} column is then appended by
          \texttt{addNewField} using \texttt{operator.mul} and a
          \texttt{".2f"} format string.

    \item The output directory (\texttt{OUTPUT\_DIR} from \texttt{.env}) is
          created if it does not exist. Both CSV files are written by
          joining headers and row values with commas.
\end{enumerate}

% ─────────────────────────────────────────────────────────────────────────────
\section{Future Improvements}

Given more time, the following changes and additions would be prioritised.

\subsection{Parameterised SQL Queries}

The current \texttt{sql\_guard} decorator checks for \texttt{;} and \texttt{--}
tokens as a defence against injection, but this approach is inherently
incomplete --- there are injection payloads that do not use these tokens. The
correct fix is to replace all f-string SQL construction with parameterised
queries using the \texttt{\%s} placeholder syntax supported by
\texttt{mysql-connector-python}:

\begin{lstlisting}[language=python, caption={Parameterised query example}]
cursor.execute(
    "SELECT * FROM %s WHERE %s = %s",
    (table, column, value)
)
\end{lstlisting}

This delegates escaping entirely to the driver and eliminates the class of
vulnerability at source.

\subsection{Replace \texttt{eval} with Explicit Type Conversion}

In \texttt{task\_three.py}, string values from the database are converted to
numbers using \texttt{eval}. While the data originates from a controlled CSV
source, \texttt{eval} is an unnecessary security risk. Replacing it with
\texttt{float()} or \texttt{int()} would be both safer and more explicit.

\subsection{A Production WSGI Server}

Flask's built-in development server is single-threaded and not suitable for
production or concurrent load testing. The application would benefit from being
served via \textbf{Gunicorn} or \textbf{uWSGI}:

\begin{lstlisting}[language=bash, caption={Running with Gunicorn}]
gunicorn -w 4 "task_two:app"
\end{lstlisting}

\subsection{Proper HTTP Status Codes and Error Handling}

The \texttt{"Customer not found"} response currently returns HTTP \texttt{200}.
A \texttt{404 Not Found} would be more semantically correct and easier to
handle by API consumers. Similarly, unhandled exceptions (e.g. database
connection failures mid-request) should be caught and returned as
\texttt{500 Internal Server Error} rather than an unhandled Python traceback.

\subsection{Database Connection Pooling}

The application opens a single persistent \texttt{mysql.connector} connection
at module import time. Under concurrent API load this will fail or produce race
conditions. Switching to \texttt{mysql.connector.pooling.MySQLConnectionPool}
or using SQLAlchemy's connection pool (already in \texttt{requirements.txt} via
\texttt{Flask-SQLAlchemy}) would resolve this.

\subsection{A Proper ORM Layer}

\texttt{Flask-SQLAlchemy} and \texttt{flask-marshmallow} are already listed as
dependencies. Migrating \texttt{utilities/database.py} to use SQLAlchemy models
would bring typed schemas, automatic migration support (via Alembic), and
eliminate the need for the hand-rolled \texttt{Database} wrapper class entirely.

\subsection{Tests}

There are currently no automated tests. A minimal test suite using
\texttt{pytest} and Flask's built-in test client would cover:

\begin{itemize}[leftmargin=2em]
    \item \texttt{GET /customer/<id>} returns correct JSON structure.
    \item \texttt{GET /customer/<id>?pretty=true} returns valid HTML.
    \item \texttt{GET /customer/INVALID} returns a \texttt{404}.
    \item \texttt{task\_three.py} output CSVs match expected headers and row counts.
\end{itemize}

\subsection{Structured Logging}

Print statements are used for status output. Replacing these with Python's
\texttt{logging} module would allow log levels to be controlled via environment
variables and log output to be directed to files or a log aggregator in
production.

\subsection{Docker Compose}

Currently the application depends on Docker being installed but manages the
container itself through \texttt{subprocess} calls. A
\texttt{docker-compose.yml} would replace this ad-hoc orchestration with a
declarative definition, making the full stack (MySQL + Python app) startable
with a single \texttt{docker compose up} command and easier to integrate into
CI pipelines.
\end{document}
